\documentclass[../main.tex]{subfiles}
\begin{document}

\chapter{Generative models}

A generative modl is a joint prob distribution $p(\boldsymbol{x})$, for $\boldsymbol{x} \in  \mathcal{X}$.

High level, there are \emph{deep generative models}(use deep nn to learn a complex mapping from a single latent vector $\boldsymbol{z}$) and \emph{probalistic graphical models} (map a set of interconnected laten vars $\boldsymbol{z}_{1}, \ldots ,\boldsymbol{z}_{L}$ to the observed variables $\boldsymbol{x}_{1}, \ldots ,\boldsymbol{x}_{D}$ using simpler mappings.) 

\begin{enumerate}
   \item variational autoencoders 
   \item autoregressive models
   \item normalizing flows
   \item diffusion models
   \item energy-based models
   \item generative adversarial networks
\end{enumerate}


\section{Goals}

\paragraph{Generating data}

we often use a \emph{conditional generative model} of the form $p(\boldsymbol{x} | \boldsymbol{c})$. $\boldsymbol{c}=$ text prompt, then it is text-to-image model.
In a discriminative model, we assume there is one correct output, wheeas in a conditional generative model, we assume there may be multiple correct outputs.

\paragraph{Density estimation}

evaluate the porb of a observed data vector, compute $p(\boldsymbol{x})$.

Use a kernel,
\[
p(\boldsymbol{x} | \mathcal{D}) = \frac{1}{N} \sum_{n=1}^{N} \mathcal{K}_{h}(\boldsymbol{x} - \boldsymbol{x}_{n}).
\]
In high dimension we need to use parametric density models $p_{\boldsymbol{\theta}}(\boldsymbol{x})$.


\paragraph{Structure discovery and Latent space interpolation}

$p(\boldsymbol{z} | \boldsymbol{x}) \propto p(\boldsymbol{z})p(\boldsymbol{x} | \boldsymbol{z})$. 


Let $\boldsymbol{x}_{1}, \boldsymbol{x}_{2}$ be two images, and let $\boldsymbol{z}_{1} = e(\boldsymbol{x}_{1})$ and $\boldsymbol{z}_{2} = e(\boldsymbol{x}_{2})$ be their latent encodings.

We can now generate new images that interpolating: $\boldsymbol{z} = \lambda \boldsymbol{z}_{1} + (1-\lambda)\boldsymbol{z}_{2}$, and then decoding by computing $\boldsymbol{x}' = d(\boldsymbol{z})$, where $d$ is the decoder.


\paragraph{Latent space arithmetic}

We can increase or decrease the amount of a disired ``semantic factor of variation''. 


\section{Evaluating generative models}

\paragraph{Likelihood-based evaluation}

To measure how close a model $q$ is to a true distribution $p$ is interms of the KL divergence
\[
D_{KL}(p || q) = \int p(\boldsymbol{x}) \log \frac{p(\boldsymbol{x})}{q(\boldsymbol{x})} = - \mathbb{H}(p)  + \mathbb{H}_{ce}(p, q).
\]
where $\mathbb{H}(p)$ is a constant and $\mathbb{H}_{ce}(p, q)$ is the cross entropy.
If $p(\boldsymbol{x})$ is a empirical distribution, the cross entropy is \emph{negative log likelihood}
\[
\text{NLL} =  \frac{1}{N} \sum_{n=1}^{N} - \log q(\boldsymbol{x}_{n}).
\]

\begin{note}
    For image and audio, the model is usually a continuous distribution $p(\boldsymbol{x}) \ge  0$ but the data is discrete. So we use \emph{uniform dequantization}, we add uniform random noise to the discrete data, and then treat it as continuous-valued data.
    This can gives a \emph{lower bound} on the average log likelihood of the discrete model on the original data.
    
    Let $\boldsymbol{z}$ be a continuous latent var, and $\boldsymbol{x}$ be a vector of binary observations computed by rounding, so $p(\boldsymbol{x} | \boldsymbol{z}) = \delta(\boldsymbol{x} - \text{round}(z))$, computed elementwise.
    We have $p(\boldsymbol{x}) = \int p(\boldsymbol{x} | \boldsymbol{z}) p(\boldsymbol{z})d\boldsymbol{z}$, let $q(\boldsymbol{z} | \boldsymbol{x})$ be prob inverse of $\boldsymbol{x}$, that is, it has support only on values where $p(\boldsymbol{x} | \boldsymbol{z}) = 1$. Jenson gives
    \[
    \begin{aligned}
    \log p(\boldsymbol{x}) \ge  \mathbb{E}_{q(\boldsymbol{z} | \boldsymbol{x})} \left[ \log p(\boldsymbol{z}) - \log q(\boldsymbol{z} | \boldsymbol{x}) \right]  \\
    \end{aligned}
    \]
    Thus if we model the density of $\boldsymbol{z} \sim  q(\boldsymbol{z} | \boldsymbol{x})$, which is a dequantized version of $\boldsymbol{x}$, we will get a lower bound on $p(\boldsymbol{x})$.

\end{note} 


\paragraph{Distances and divergences in feature space}

We want to compare distributions in a semantically meaningful way.

The \emph{Inception score} measures the average KL divergence between the marginal distribution of class labels obtaiined from samples 
\[
p_{\boldsymbol{\theta}}(y) = \int p_{\text{disc}} (y | \boldsymbol{x}) p_{\boldsymbol{\theta}}(\boldsymbol{x}) d\boldsymbol{x}
\]
This leads to the following score:
\[
\text{IS} = \exp \left[ \mathbb{E}_{p_{\boldsymbol{\theta}(x)}} D_{\mathbb{KL}} \left( p_{\text{disc}}(Y|\boldsymbol{x}) || p_{\boldsymbol{\theta}(Y)} \right)  \right] 
\]

\paragraph{Precision and recall metrics}


\end{document}
